\documentclass[10pt,a4paper]{amsart}
\usepackage[margin=19mm]{geometry}
\usepackage{amsmath,amssymb,mathtools}
\usepackage[scr=rsfs,cal=euler]{mathalfa}
\usepackage{xfrac}
\usepackage[unicode,hidelinks,bookmarks=false]{hyperref}
\newcommand{\ie}{i.e.\ }
\newcommand{\cf}{cf.\ }
\newcommand{\resp}{respectively\ }
\newcommand{\Subsection}{\subsection}
\providecommand{\nobreakdash}{\nobreak\hbox{-}}
\setlength{\emergencystretch}{2em}
\multlinegap0pt
\usepackage{csquotes}
\usepackage{amssymb}
\usepackage{mathtools}
\usepackage[shortlabels]{enumitem}
\usepackage{xcolor}
\newtheorem{theorem}{Theorem}
\title{The Schwarzian Derivative for Convex Holomorphic Mappings in Several Complex Variables}
\author{Rodrigo Hern\'andez}
\date{Source: arXiv:2604.12010v1 (CC BY 4.0)}
\begin{document}
\maketitle

\noindent\emph{Source and licence.} The Schwarzian Derivative for Convex Holomorphic Mappings in Several Complex Variables. Version arXiv:2604.12010v1 (CC BY 4.0), distributed by arXiv under the Creative Commons Attribution 4.0 International licence, http://creativecommons.org/licenses/by/4.0/, as recorded in the arXiv licence metadata for this exact version and shown on the abstract page https://arxiv.org/abs/2604.12010v1. Full licence notice: \texttt{licenses/arxiv-2604.12010-CC-BY-4.0.txt}.\par\smallskip

\noindent\textit{Editorial scope.} Counted unit: the article's theorem theorem (source0.tex lines 383--390). The article's complete original proof of it is delivered with it (source0.tex lines 392--461, one \texttt{proof} environment of 70 lines). No mathematics is written, repaired or re-derived: the statement and the proof are the article's own lines, in the article's own notation, and no retained line is substituted or re-flowed.  No further source-local context is needed: every cross-reference of every retained line resolves inside the two delivered spans.  Nothing was omitted from any retained span, so the package carries no omission record.  No retained line contains a citation, so the wrapper carries no bibliography and no bibliography entry.  No retained line contains a figure environment, an image-inclusion command or drawing code, so no image is delivered and the package declares no asset dependency.  Editorial formatting only, in addition: an AMS \texttt{amsart} preamble that restates the article's own declarations and macros used by the retained lines, namely the environments \texttt{theorem} on separate counters, together with the standard packages those lines need.  The retained environments print in this document as Theorem 1 in order of appearance, whereas the article prints the same items with the numbering of its own full document.  The complete native source of the article as distributed in the arXiv e-print for this exact version is bundled unchanged as \texttt{sources/198376/source0.tex}.
\begin{theorem}
Let $f:\mathbb{P}^n\to\mathbb{C}^n$ be a locally univalent convex mapping.
Then
\[
\|S_f\|\leq \frac{\sqrt{8(n+3)(n-1)}}{n+1}.
\]
The bound is sharp.
\end{theorem}

\begin{proof}
Let $\vec{v}$ be a unit vector in the Bergman metric of the polydisk, so that
\[
\|\vec{v}\|=\left[2\sum_{i=1}^n\frac{|v_i|^2}{(1-|z_i|^2)^2}\right]^{1/2}=1.
\]
Then, since $S_f(z)(\vec{v},\vec{v})=(\delta_k v_k)_k$ and the Bergman norm
of this vector is bounded by $\max_k|\delta_k|\cdot\|\vec{v}\|$, we have
$\|S_f(z)\|\leq m(z)\|\vec{v}\|=m(z)$, where
$m(z)=\max\{|m_k(z)|:k=1,\ldots,n\}$ and
\[
m_k(z) =\left(\frac{-2}{n+1}\frac{\varphi_1''}{\varphi_1'}(z_1)v_1
\right)
+\cdots+\left(\frac{n-1}{n+1}\frac{\varphi_k''}{\varphi_k'}(z_k)v_k
\right)
+\cdots+\left(\frac{-2}{n+1}\frac{\varphi_n''}{\varphi_n'}(z_n)v_n
\right).
\]
Rewriting $m_k$ as
\[
|m_k(z)|=
\left|a_1\cdot\frac{v_1}{1-|z_1|^2}+\cdots+a_n\frac{v_n}{1-|z_n|^2}\right|,
\]
where
\[
a_i=\frac{-2}{n+1}\frac{\varphi_i''}{\varphi_i'}(z_i)(1-|z_i|^2),
\quad i\neq k,
\qquad
a_k=\frac{n-1}{n+1}\frac{\varphi_k''}{\varphi_k'}(z_k)(1-|z_k|^2),
\]
and applying the Cauchy--Schwarz inequality, we obtain
\[
|m_k(z)| \leq
\left(\sum_{i=1}^{n}|a_i|^2\right)^{1/2}
\left(\sum_{i=1}^{n}\frac{|v_i|^2}{(1-|z_i|^2)^2}\right)^{1/2}
=\left(\sum_{i=1}^{n}|a_i|^2\right)^{1/2}\frac{1}{\sqrt{2}}.
\]
Since each mapping $\varphi_i$ is convex, the standard estimate
$|\varphi_i''/\varphi_i'(z_i)|(1-|z_i|^2)\leq 4$ holds, which implies that
\[
\left(\sum_{i=1}^{n}|a_i|^2\right)^{1/2}
\leq
\left(\frac{4}{(n+1)^2}\cdot16(n-1)+16\left(\frac{n-1}{n+1}\right)^2\right)^{1/2}
=\left(\frac{16(n-1)(n+3)}{(n+1)^2}\right)^{1/2}.
\]
This yields
\[
\|S_f\|\leq \frac{2\sqrt{2(n+3)(n-1)}}{n+1}.
\]

Sharpness is achieved for $f=(\varphi_1,\ldots,\varphi_n)$ with
$\varphi_i(z_i)=1/(1-z_i)$ for all $i=1,\ldots,n$. In this case,
\[
|a_i|=\left|\frac{4(1-|z_i|^2)}{(n+1)(1-z_i)}\right|\leq \frac{4(1+|z_i|)}{n+1},
\qquad
a_k=\frac{2(n-1)(1-|z_k|^2)}{(n+1)(1-z_k)},
\]
and equality holds when $z_i\in[0,1)$; moreover, $|a_i|\to 8/(n+1)$ and
$|a_k|\to 4(n-1)/(n+1)$ as $z_i$ and $z_k$ tend to $1$ along the real axis.
For any $z\in\mathbb{P}^n$ one may take
\[
v_i=\frac{\overline{a_i}(1-|z_i|^2)}{\sqrt{2(|a_1|^2+\cdots+|a_n|^2)}},
\quad i=1,\ldots,n,
\]
so that
\[
\lim_{z_i\to 1^-}m_k(z)
=\lim_{z_i\to 1^-}\frac{\sqrt{|a_1|^2+\cdots+|a_n|^2}}{\sqrt{2}}
= \frac{2\sqrt{2(n+3)(n-1)}}{n+1}. \qedhere
\]
\end{proof}

\end{document}
